%% Chaîne énergétique : batterie -> moteur électrique -> disque en rotation.
%% Règle de lecture : les formes d'énergie sont encadrées, le convertisseur est
%% entouré (ellipse), la nature du transfert est écrite au-dessus des flèches.
\documentclass[tikz,border=4pt]{standalone}
\usepackage{liaisons}
\usetikzlibrary{shapes.geometric, positioning}
\begin{document}
\begin{tikzpicture}[x=1cm,y=1cm,
  forme/.style={draw=solideB, line width=1pt, align=center, minimum width=1.5cm,
                minimum height=1.05cm, inner sep=2pt,
                font=\scriptsize\bfseries, text=solideB},
  conv/.style={ellipse, draw=solideE, line width=1pt, align=center,
               minimum width=1.8cm, minimum height=1.1cm, inner sep=1pt,
               font=\scriptsize\bfseries, text=solideE},
  titre/.style={font=\scriptsize, text=black!70, align=center},
  fleche/.style={-{Stealth[length=7pt]}, line width=1.2pt, draw=solideE},
  etiq/.style={font=\scriptsize\bfseries, text=solideD, inner sep=2pt}]

%% ---- les trois maillons de la chaîne -------------------------------------
\node[forme] (chim) at (0,0)   {Énergie\\chimique};
\node[conv]  (mot)  at (3.15,0) {Moteur\\électrique};
\node[forme] (cin)  at (6.3,0) {Énergie\\cinétique};

%% ---- titres au-dessus de chaque maillon ----------------------------------
\node[titre, above=3pt of chim] {Batterie};
\node[titre, above=3pt of mot]  {Convertisseur};
\node[titre, above=3pt of cin]  {Disque\\en rotation};

%% ---- transferts ----------------------------------------------------------
\draw[fleche] (chim) -- node[etiq, above] {Électrique} (mot);
\draw[fleche] (mot)  -- node[etiq, above] {Mécanique}  (cin);

\node[font=\scriptsize, text=black!70] (env) at (5.0,-1.7) {Environnement};
\draw[fleche] (mot.south) -- node[etiq, below left=-1pt] {Thermique} (env.north west);
\end{tikzpicture}
\end{document}
